\chapter{Native Python Tasks}
\label{ch:native-tasks}

\section{Design Goals}

\texttt{TaskType.NATIVE} lets an \emph{importable Python function} act as a
workflow task while preserving DAGonStar dependencies, scratch isolation,
staged files, declared outputs, and checkpoint metadata. Native tasks are
appropriate when scientific logic is naturally Python code --- for instance,
an array-processing or statistical routine built on NumPy or xarray --- and
should not be artificially encoded as a shell command merely to fit the
\texttt{BATCH} task model.

\begin{cautionbox}
\texttt{NativeTask} records a \texttt{module:function} reference and
JSON-compatible parameter bindings. The runner \emph{imports} the function in
the task environment, so functions must be module-level and importable.
Lambdas and nested functions are deliberately rejected, because they are not
portable across executors (in particular, a remote Slurm executor must be
able to import the same function from the same module on a different
machine).
\end{cautionbox}

\section{Worked Example}

\begin{lstlisting}[caption={\texttt{native\_functions.py}: the importable native function.}]
def transform(input_file: str, scale: float, output_file: str) -> dict:
    with open(input_file, encoding="utf-8") as source:
        values = [float(line.strip()) for line in source if line.strip()]
    with open(output_file, "w", encoding="utf-8") as destination:
        for value in values:
            destination.write(f"{value * scale}\n")
    return {"count": len(values), "scale": scale}
\end{lstlisting}

\begin{lstlisting}[caption={\texttt{native\_workflow.py}: the workflow that calls it.}]
import json
import tempfile
from pathlib import Path

from dagon import Workflow
from dagon.task import DagonTask, TaskType


def config(scratch):
    return {
        "batch": {"scratch_dir_base": scratch, "run_base": "", "threads": "1"},
        "dagon_service": {"use": "False", "route": "http://localhost:57000"},
        "ftp_pub": {"ip": "localhost"},
        "slurm": {"partition": ""},
    }


with tempfile.TemporaryDirectory(prefix="dagon-native-lesson-") as scratch:
    workflow = Workflow("Lesson14", config=config(scratch))
    produce = DagonTask(
        TaskType.BATCH, "produce", "mkdir -p data; printf '2\n4\n' > data/input.txt"
    )
    transform = DagonTask(
        TaskType.NATIVE,
        "transform",
        "native_functions:transform",
        inputs={"input_file": "workflow:///produce/data/input.txt", "scale": 1.5},
        outputs={"output_file": "scaled.txt"},
    )
    consume = DagonTask(
        TaskType.BATCH, "consume", "cat workflow:///transform/outputs/scaled.txt"
    )
    for task in (produce, transform, consume):
        workflow.add_task(task)
    workflow.run()
    print(Path(transform.working_dir, "outputs", "scaled.txt").read_text())
    print(json.loads(Path(transform.working_dir, ".dagon", "native_result.json").read_text()))
\end{lstlisting}

\begin{practicebox}
This small two-file layout (\texttt{native\_functions.py} alongside the
workflow script) is essential, not incidental: it is precisely what makes the
function importable by the separate runner process. A native function
defined inline inside the workflow script's \texttt{\_\_main\_\_} block would
not be importable, and would fail at the same architectural boundary that
excludes lambdas.
\end{practicebox}

\begin{figure}[h]
\centering
\begin{tikzpicture}[
  node distance=20mm,
  task/.style={rectangle, rounded corners, draw=dagonblue, very thick,
    fill=blue!4, minimum width=30mm, minimum height=13mm, font=\sffamily\small, align=center},
  arr/.style={-{Stealth[length=2.6mm]}, thick, dagonblue}
]
  \node[task] (prod) {produce\\ \footnotesize \texttt{data/input.txt}};
  \node[task, right=26mm of prod] (nat) {transform\\ \footnotesize \texttt{TaskType.NATIVE}\\ \texttt{native\_functions:transform}};
  \node[task, right=26mm of nat] (cons) {consume\\ \footnotesize reads \texttt{outputs/scaled.txt}};
  \draw[arr] (prod) -- node[above,font=\sffamily\scriptsize]{\texttt{inputs=\{input\_file: workflow://...\}}} (nat);
  \draw[arr] (nat) -- node[above,font=\sffamily\scriptsize]{\texttt{outputs=\{output\_file: scaled.txt\}}} (cons);
\end{tikzpicture}
\caption{The Native task pipeline: a staged input file becomes a Python
function argument, and a declared output parameter becomes an ordinary,
dependency-trackable file for the downstream consumer.}
\label{fig:native-task-flow}
\end{figure}

\section{Interpretation}

The reference creates \texttt{produce} $\to$ \texttt{transform}; the runner
passes a scratch-local file path to the \texttt{input\_file} parameter and a
scratch-local output path to the \texttt{output\_file} parameter. The
declared output is therefore available to the downstream task through
\texttt{workflow:///transform/outputs/scaled.txt}, exactly as any other
task's file output would be. The Python function's \emph{return value} is
retained as JSON metadata in \texttt{.dagon/native\_result.json} ---
analogous in spirit to the full-response retention strategy used by LLM
tasks (\Cref{ch:llm-tasks}): both preserve a richer, machine-readable record
than the minimum a human reader strictly needs.

\section{Constraints and Extension}

\begin{itemize}[leftmargin=1.4em]
  \item Scalars and return values must be JSON-serializable.
  \item Output paths must be relative and cannot escape the task's scratch
    directory.
  \item For a configured cluster, \texttt{executor="slurm"} may be used with
    scheduler settings supplied in \texttt{resources}; the function's module
    and DAGonStar itself must be installed where the compute node can import
    them.
\end{itemize}

\section{Verification}

\begin{lstlisting}[style=bash]
python3 -m unittest tests.test_native -v
\end{lstlisting}

The course examples \texttt{examples/18-native/native\_task.py} and
\texttt{examples/19-native-bindings/native\_bindings.py} (Lectures~18--19)
exercise the same task type with binding validation and Slurm resources.
See also \texttt{examples/native/README.md} in the DAGonStar repository for further,
runnable verification material.

\section{Native Tasks in a Scientific Pipeline}

Native tasks complete a triad with LLM and Web tasks that, taken together,
constitute the engine's full AI--HPC integration surface:

\begin{table}[h]
\centering
\caption{The three non-shell task types and their characteristic role.}
\label{tab:ai-hpc-triad}
\begin{tabular}{@{}lll@{}}
\toprule
\textbf{Task type} & \textbf{What it executes} & \textbf{Characteristic role} \\
\midrule
\texttt{LLM}    & OpenAI-compatible chat request & AI-assisted interpretation, summarization \\
\texttt{WEB}    & Arbitrary HTTP/HTTPS request   & Publication, ingestion, third-party API calls \\
\texttt{NATIVE} & Importable Python function     & In-process scientific computation (NumPy/xarray-style) \\
\bottomrule
\end{tabular}
\end{table}

In a typical environmental forecasting chain (\Cref{ch:case-study}), a
\texttt{NATIVE} task is often the most natural choice for a post-processing
step that is already implemented as a Python function --- for instance, a
unit conversion, a regridding routine, or a statistical quality-control
check --- avoiding the indirection of wrapping that function in a throwaway
command-line script merely to satisfy a \texttt{BATCH} task interface.

\section{Why Importability, Not Mere Callability, Is Required}

\Cref{ch:native-tasks} stated, as a constraint, that native functions must
be module-level and importable, and that lambdas and nested functions are
rejected. It is worth explaining \emph{why} this constraint is structurally
necessary, rather than an arbitrary restriction, since understanding the
reason clarifies what a workflow author can and cannot work around.

A \texttt{NativeTask} does not execute the supplied function inside the same
Python process, or even necessarily the same thread, that constructed the
\texttt{DagonTask} object. Recall from \Cref{sec:threads-not-processes} that
DAGonStar tasks are threads, not processes, for backends whose computation
happens in a spawned subprocess or over a network connection; a
\texttt{NativeTask} configured with \texttt{executor="local"} does still
ultimately need to make the target function available to whatever runner
process or thread actually performs the call, and a \texttt{NativeTask}
configured with \texttt{executor="slurm"} must additionally make the
function available on a \emph{different machine entirely} --- the compute
node a Slurm job lands on, which may share no process memory whatsoever
with the controller that submitted the job.

A lambda or a function nested inside another function exists only as an
in-memory object inside the defining process; Python provides no general
mechanism to serialize such an object's code and ship it, intact, to another
process or machine (the standard \texttt{pickle} mechanism, which many
readers may expect to handle this, explicitly does not support pickling
lambdas or local functions, for exactly this reason). A module-level
function, by contrast, can always be \emph{re-imported by name} --- the
runner simply executes the equivalent of \texttt{from native\_functions
import transform} --- which works identically whether the runner is in the
same process, a different local process, or a remote compute node, provided
only that the same module is installed and importable wherever the runner
executes.

\begin{practicebox}
This is the same architectural principle already encountered, in a
different guise, in this book's discussion of remote Batch and Slurm
execution (\Cref{ch:backends}): just as a \texttt{BATCH} task's portability
across hosts depends on the target executable being installed wherever the
command runs, a \texttt{NATIVE} task's portability across executors depends
on the target Python module being installed wherever the runner executes.
The module-level-function constraint is simply this same requirement,
restated for Python's own import mechanism rather than for a shell's
\texttt{PATH}.
\end{practicebox}

\section{Native Tasks and the Scientific Python Ecosystem}

\Cref{ch:native-tasks} already suggested that Native tasks are a natural fit
for NumPy- or xarray-style computation; it is worth developing this
suggestion into a fuller discussion of where, exactly, the boundary should
fall between a \texttt{NATIVE} task and a \texttt{BATCH} task that happens
to invoke a Python script.

Consider a regridding or unit-conversion routine already implemented as a
well-tested, importable function in a research group's existing Python
codebase. Wrapping it as a \texttt{BATCH} task would require writing a
small command-line entry point --- parsing arguments, reading and writing
files explicitly --- purely to satisfy the shell-command interface, even
though the underlying logic never needed a command-line interface in the
first place. Wrapping the same function as a \texttt{NATIVE} task, by
contrast, requires no such adapter: the function's existing signature,
provided its arguments are JSON-compatible scalars or file paths, can be
bound directly through the \texttt{inputs}/\texttt{outputs} mapping
mechanism of Lecture~18, with DAGonStar itself handling the file staging and
result capture that a hand-written command-line wrapper would otherwise
need to implement redundantly.

\begin{example}
Contrast the two wrapping strategies for a hypothetical
\texttt{convert\_units(value, from\_unit, to\_unit)} function already
present in a research codebase. As a \texttt{BATCH} task, a researcher would
need to write and maintain a separate \texttt{convert\_units\_cli.py}
script with \texttt{argparse}-based argument handling, purely to make the
function invocable from a shell command string. As a \texttt{NATIVE} task,
the existing function can be called directly --- provided it is moved to,
or already lives in, an importable module --- via:
\begin{lstlisting}
convert = DagonTask(
    TaskType.NATIVE, "convert", "units_module:convert_units",
    inputs={"value": 18.5, "from_unit": "celsius", "to_unit": "kelvin"},
    outputs={},
)
\end{lstlisting}
eliminating the redundant command-line adapter entirely, at the cost of the
importability constraint discussed above.
\end{example}

\begin{cautionbox}
This convenience is not unconditional. A function that internally launches
a separate, long-running compiled simulation via its own subprocess calls
gains little from \texttt{NATIVE}-task wrapping over a direct
\texttt{BATCH} or \texttt{SLURM} task targeting the simulation executable
itself; the JSON-serializability requirement on inputs and outputs
(\Cref{ch:native-tasks}) is also a genuine constraint that purely in-process,
array-oriented scientific computation satisfies easily, but that functions
exchanging large binary scientific data structures --- raw NumPy arrays, for
instance, as opposed to file paths to NetCDF files --- do not satisfy
directly without an explicit serialization step the function itself must
perform.
\end{cautionbox}

\section{Chapter Summary}

\texttt{TaskType.NATIVE} closes Part~II's progression through heterogeneous
execution backends by bringing ordinary, importable Python functions into
the same dependency, staging, and provenance model already established for
shell commands (Part~I), containers and HPC schedulers
(\Cref{ch:backends}), and AI/web services (\Cref{ch:llm-tasks,ch:web-tasks}).
This chapter additionally explained, in terms of Python's own serialization
limitations, why the module-level-importability constraint is structurally
necessary rather than arbitrary, and examined where the boundary should fall
between wrapping existing scientific Python code as a \texttt{NATIVE} task
versus as a \texttt{BATCH} task invoking a hand-written command-line
adapter. \Cref{ch:external-software} now turns from the engine's task-type
taxonomy to the practical discipline of wrapping real, compiled scientific
software --- the discipline that the environmental-science case study
depends on directly.

\begin{exercise}
A colleague proposes converting a research group's existing
\texttt{archive2plot.py} rendering script (\Cref{ch:case-study}) from a
\texttt{BATCH} task, invoked as a standalone command-line program, into a
\texttt{NATIVE} task calling its core rendering function directly. List the
specific, concrete checks you would perform on the existing script before
agreeing this conversion is straightforward, referencing both the
importability constraint and the JSON-serializability constraint discussed
in this chapter.
\end{exercise}
